\documentclass[a4paper,10pt,landscape]{article}

\usepackage[czech]{babel}
\usepackage[T1]{fontenc}
\usepackage[utf8]{inputenc}
\usepackage{amsmath,amssymb}
\usepackage{geometry}
\usepackage{enumitem}

\geometry{a4paper,landscape,left=8mm,right=8mm,top=8mm,bottom=8mm}
\pagestyle{empty}
\setlength{\parindent}{0pt}
\setlength{\parskip}{0pt}
\setlist[enumerate]{leftmargin=*,itemsep=2pt,topsep=2pt}

% Sazba podle sablony/testy/A5x2.tex, beze změny společné šablony.
% Zdroj zadání: přímý pokyn učitele v konverzaci 27. 10. 2026,
% zpráva „Vytvoř písemný test z matematiky na téma aritmetická posloupnost“.
% Učitelem uvedená banka: „Aritmetická posloupnost – 100 úloh“.
% A1–A5: úlohy banky 22, 27, 31, 54, 64.
% B1–B5: úlohy banky 23, 28, 32, 55, 65.
% Banka nebyla v projektu nalezena; čísla jsou evidována podle učitele,
% nebyla nezávisle porovnána s původním souborem banky. Přesným podkladem
% pro znění a hodnoty je úplné zadání učitele v uvedené zprávě.
% Starší verze tohoto konkrétního testu ani její bodování nebyly nalezeny.
% import_testu/bezne/testPosloupnosti.tex obsahuje jiný test (3 úlohy).
% Zařazení: 4_rocnik/tematicke_plany/SPS.pdf, s. 1, body 1.1–1.2.
\newcommand{\TestTitle}{Aritmetická posloupnost}
\newcommand{\TestHeader}[1]{%
    {\bfseries TEST -- \TestTitle}\hfill{\bfseries Varianta #1}\\[2mm]
    Jméno a příjmení: \rule{4.5cm}{0.4pt}\hfill Třída: \rule{1.5cm}{0.4pt}\\[2mm]
    Datum: \rule{2.4cm}{0.4pt}\hfill Počet bodů: \rule{1.4cm}{0.4pt}\hfill Známka: \rule{1.4cm}{0.4pt}\\[1mm]
    \hrule\vspace{3mm}
}
\newcommand{\AnswerSpace}[1]{\vspace{#1}}

\begin{document}
\noindent
\begin{minipage}[t]{0.48\textwidth}
    \small
    \TestHeader{A}

    % ZADÁNÍ VARIANTY A
    \begin{enumerate}
        \item Jsou dány členy $a_4=6$, $a_{11}=34$.
        Určete $d$, $a_1$ a $a_8$.
        \AnswerSpace{26mm}

        \item Určete $n\in\mathbb{N}$, pro které platí $a_n=77$,
        jestliže $a_1=2$, $d=5$.
        \AnswerSpace{22mm}

        \item Najděte vzorec pro $n$-tý člen posloupnosti
        $7,\ 11,\ 15,\ 19,\ldots$
        \AnswerSpace{21mm}

        \item Určete počet členů posloupnosti $4,\ 9,\ 14,\ldots,149$
        a jejich součet.
        \AnswerSpace{27mm}

        \item Mezi čísla $2$ a $26$ vložte pět čísel tak, aby vzniklo sedm
        po sobě jdoucích členů aritmetické posloupnosti.
        \AnswerSpace{24mm}
    \end{enumerate}
\end{minipage}%
\hfill\vrule width 0.4pt\hfill
\begin{minipage}[t]{0.48\textwidth}
    \small
    \TestHeader{B}

    % ZADÁNÍ VARIANTY B
    \begin{enumerate}
        \item Jsou dány členy $a_7=2$, $a_{10}=-2$.
        Určete $d$, $a_1$ a $a_{20}$.
        \AnswerSpace{26mm}

        \item Určete $n\in\mathbb{N}$, pro které platí $a_n=-53$,
        jestliže $a_1=11$, $d=-4$.
        \AnswerSpace{22mm}

        \item Najděte vzorec pro $n$-tý člen posloupnosti
        $12,\ 7,\ 2,\ {-3},\ldots$
        \AnswerSpace{21mm}

        \item Určete počet členů posloupnosti $120,\ 113,\ 106,\ldots,-41$
        a jejich součet.
        \AnswerSpace{27mm}

        \item Mezi čísla $-8$ a $34$ vložte šest čísel tak, aby vznikla
        aritmetická posloupnost.
        \AnswerSpace{24mm}
    \end{enumerate}
\end{minipage}
\end{document}
